% Generated mechanically from frozen input p12/ja/build/ch113/desirables_full.tex.
% Do not edit this generated file; edit the bound locale source instead.

% OK, start here.
%

\setcounter{chapter}{112}
\chapter{要望事項}


\phantomsection
\label{desirables-section-phantom}



\section{序論}
\label{desirables-section-introduction}

\noindent
これは基本的に，Stacks プロジェクトに盛り込みたい事項を列挙しただけのものである．
Stacks プロジェクトには継続的に内容を追加しているため，この一覧は常に
その時点でのプロジェクトの状態より多少遅れている．実際，追加すべき事項を
列挙しようとしたこと自体が誤りだったのかもしれない．この一覧を最新に保つのは
不可能に思われるからである．

\medskip\noindent
最終更新：2017 年 8 月 31 日（木）．


\section{規約}
\label{desirables-section-conventions}

\noindent
この文書で用いる規約を短く列挙した章が必要である．その章はすでに存在する；
規約，節 \ref{conventions-section-comments} を参照せよ．
しかし，そこにはさらに多くのことを追加できる．とりわけ「隠れた」規約や
暗黙の仮定を見つけ，そこに記録することが有用であろう．


\section{サイトとトポス}
\label{desirables-section-sites}

\noindent
サイトと層に関する章がある；
サイトと層，節 \ref{sites-section-introduction} を参照せよ．
環付きサイト（およびトポス）とその上の加群に関する章もある；
サイト上の加群，節 \ref{sites-modules-section-introduction} を参照せよ．
さらに，この設定でのコホモロジーに関する章もある；
サイト上のコホモロジー，節 \ref{sites-cohomology-section-introduction} を参照せよ．
しかし，とりわけコホモロジーの章には，なお多くの内容を追加できる．


\section{スタック}
\label{desirables-section-stacks}

\noindent
（抽象的な）スタックに関する章がある；
スタック，節 \ref{stacks-section-introduction} を参照せよ．
次のことが実現すれば望ましい：
\begin{enumerate}
\item 「スタック化」の議論を改善する，
\item スタック化の例を与える，
\item 一般に，さらに多くの例を加える，
\item gerbe の議論を改善する．
\end{enumerate}
まだ追加されていない結果の例を挙げる．$\mathcal{C}$ 上の可換群の層
$\mathcal{F}$ が与えられたとき，$\mathcal{F}$ を帯とする gerbe の
同値類全体の集合は $H^2(\mathcal{C}, \mathcal{F})$ と全単射に対応する．


\section{単体的方法}
\label{desirables-section-simplicial}

\noindent
単体的方法に関する章がある；
単体的方法，節 \ref{simplicial-section-introduction} を参照せよ．
この章は見直して改善する必要がある．単体的ホモトピー（組合せ的ホモトピーともいう）と
Kan 複体との関係についての議論を改善すべきである．
単体的代数空間に関する章もある；
単体的代数空間，節 \ref{spaces-simplicial-section-introduction} を参照せよ．
この章では，単体的位相空間，単体的サイト，単体的トポスを簡単に論じている．
さらに「単体的代数幾何」を展開し，単体的スキーム（あるいは単体的代数空間，
単体的代数スタック）を論じ，幾何学的問題やそれらのコホモロジーなどを扱うことができる．


\section{スキームのコホモロジー}
\label{desirables-section-schemes-cohomology}

\noindent
準連接層のコホモロジーに関する章はすでにある；
スキームのコホモロジー，節 \ref{coherent-section-introduction} を参照せよ．
スキーム上の準連接層の導来圏を論じる章もある；
スキームの導来圏，節 \ref{perfect-section-introduction} を参照せよ．
さらに，Noether スキームの双対性とスキームの射の相対双対性を論じる章もある；
スキームの双対性，節 \ref{duality-section-introduction} を参照せよ．
スキームの \'etale コホモロジーとクリスタリンコホモロジーに関する章もある．
しかし，これらの章の内容の大半は非常に基礎的であり，なお多くのことを
追加できるし，追加すべきである．

\section{Schlessinger 流の変形理論}
\label{desirables-section-deformation-schlessinger}

\noindent
この内容に関する章がある；
形式変形理論，節 \ref{formal-defos-section-introduction} を参照せよ．
一般論の例を論じる章もある；
変形問題の例，節 \ref{examples-defos-section-introduction} を参照せよ．
さらに，変形理論，節 \ref{defos-section-introduction} では，
環（および加群）の変形，環付き空間（および加群の層）の変形，
環付きトポス（および加群の層）の変形を論じる．
この章では，素朴な余接複体を用いて，障害，一次変形，および無限小自己同型を記述する．
この内容は，後の章でモジュライスタックの代数性を示す際にいくつか応用されている．
完全な余接複体を論じる章もある；
余接複体，節 \ref{cotangent-section-introduction} を参照せよ．


\section{代数スタックの定義}
\label{desirables-section-definition-algebraic-stacks}

\noindent
代数スタックとは，fppf 位相を入れたスキームの圏上の亜群のスタックであって，
その対角射が代数空間によって表現可能であり，かつあるスキームからの
全射的滑らかな射の終域となるものである．
代数スタック，節 \ref{algebraic-section-algebraic-stacks} を参照せよ．
「Deligne--Mumford スタック」とは，Deligne と Mumford の論文 \cite{DM} におけるように，
あるスキームからそれへの全射的 \'etale 射が存在する代数スタックである；
代数スタック，定義 \ref{algebraic-definition-deligne-mumford} を参照せよ．
「Artin スタック」という語は，Artin の論文に現れる種類のスタックのために留保する；
\cite{ArtinI}，\cite{ArtinII}，および \cite{ArtinVersal} を参照せよ．
一つの可能な定義は，局所 Noether スキーム $S$ 上の代数スタック
$\mathcal{X}$ であって，$\mathcal{X} \to S$ が局所有限型であるものを
Artin スタックとすることである\footnote{すなわち，これらはちょうど，
$S$ 上の代数スタックであって，Artin の公理，節 \ref{artin-section-axioms} の
Artin の公理 [-1], [0], [1], [2], [3], [4], [5] を満たすものである．}．


\section{スキーム，代数空間，代数スタックの例}
\label{desirables-section-examples-stacks}

\noindent
Stacks プロジェクトには現在，モジュライスタックとその性質を論じる章が二つある；
モジュライスタック，節 \ref{moduli-section-introduction} および
曲線のモジュライ，節 \ref{moduli-curves-section-introduction} を参照せよ．
今後，たとえば次のものを追加する予定である：
\begin{enumerate}
\item $\mathcal{A}_g$，すなわち主偏極をもつ種数 $g$ のアーベルスキーム，
\item $\mathcal{A}_1 = \mathcal{M}_{1, 1}$，すなわち $1$ 点付きの滑らかな
射影的種数 $1$ 曲線，
\item $\mathcal{M}_{g, n}$，すなわち互いに異なるラベル付きの $n$ 点をもつ
滑らかな射影的種数 $g$ 曲線，
\item $\overline{\mathcal{M}}_{g, n}$，すなわち安定な $n$ 点付き節点射影的種数 $g$ 曲線，
\item $\SheafHom_S(\mathcal{X}, \mathcal{Y})$，すなわち射のモジュライ
（スタック $\mathcal{X}$，$\mathcal{Y}$ と基底スキーム $S$ に適切な条件を課す），
\item $\textit{Bun}_G(X) = \SheafHom_S(X, BG)$，すなわち幾何学的 Langlands 計画に現れる
$G$-束のスタック（スキーム $X$，群スキーム $G$，基底スキーム $S$ に適切な条件を課す），
\item $\Picardstack_{\mathcal{X}/S}$，すなわち基底スキーム（または空間）上の
代数スタックに付随する Picard スタック．
\end{enumerate}
より一般に，Stacks プロジェクトには幾何学的に意味のある例がやや不足している．


\section{代数スタックの性質}
\label{desirables-section-stacks-properties}

\noindent
これはおそらく，比較的取り組みやすいプロジェクトの一つである．
基礎理論の大部分はすでに整っているからである．もちろん，これらは実際には
スタックの射の性質である．特異点を（滑らかな因子を除いて）定義することなどができる．
連結な正規スタックが既約であることなどを証明せよ．

\section{代数スタックの lisse-\'etale サイト}
\label{desirables-section-lisse-etale}

\noindent
これはすでに，代数スタックのコホモロジー，節
\ref{stacks-cohomology-section-lisse-etale} で導入されている．
これが代数スタックの $1$-射に関して関手的でないことを示す例は，
例，節 \ref{examples-section-lisse-etale-not-functorial} で論じられている．
もちろん，これについてはさらに多くのことを述べられるが，できる限り
「大きな」\'etale サイトを用いて命題を証明することが非常に有用だと分かっている．



\section{いつも知りたかったが尋ねるのが怖かったこと}
\label{desirables-section-stacks-fun-lemmas}

\noindent
繰り返し用いる有用な補題のうち，文献に明示的には述べられていないもの，
あるいは参照先を見つけにくいものが数多くあるはずである．いわば手法集である．

\medskip\noindent
例：スキームにおける二つの亜群 $R\Rightarrow U$ と
$R' \Rightarrow U'$ が与えられたとき，$1$-射
$[U/R] \to [U'/R']$ をもつということは，スキームにおける亜群だけを用いて
どのように記述されるか．



\section{スタック上の準連接層}
\label{desirables-section-quasi-coherent}

\noindent
これらは，代数スタックのコホモロジー，節
\ref{stacks-cohomology-section-introduction} の章で定義され，論じられている．
加群の導来圏は，スタックの導来圏，節
\ref{stacks-perfect-section-introduction} の章で論じられている．
これらの章には，さらに多くの内容を追加できる．



\section{平坦性と滑らかさ}
\label{desirables-section-flat-smooth}

\noindent
スキームからの平坦全射をもつという条件が滑らかな全射をもつという条件の
代わりになる，という Artin の定理である．これは現在，表現可能性の判定条件，定理
\ref{criteria-theorem-bootstrap} として利用できる．


\section{Artin の表現可能性定理}
\label{desirables-section-representability}

\noindent
これは Artin の公理，節 \ref{artin-section-introduction} の章で論じられている．
応用も一つある；Quot 空間と Hilbert 空間，定理
\ref{quot-theorem-coherent-algebraic} を参照せよ．
さらに多くの応用を加えるべきであり，章そのものも整理する必要がある．


\section{DM スタックはスキームによって有限被覆される}
\label{desirables-section-dm-finite-cover}

\noindent
代数空間についての対応する結果はすでにある；
代数空間の極限，節 \ref{spaces-limits-section-finite-cover} を参照せよ．
不足しているのは，DM スタックおよび準 DM スタックに対する結果である．


\section{固有性に関する Martin Olsson の論文}
\label{desirables-section-proper-parametrization}

\noindent
この論文は，固有性の二つの概念が同じであることを証明する．その前半は現在，
代数スタックに対する Chow の補題という形で利用できる；
スタックの射の補足，定理
\ref{stacks-more-morphisms-theorem-chow-finite-type} を参照せよ．
その帰結として，ある場合には代数スタックの固有性の付値判定法を確認する際，
DVR のみを用いれば十分であることを示す；スタックの射の補足，節
\ref{stacks-more-morphisms-section-Noetherian-valuative-criterion} を参照せよ．


\section{連接層の固有順像}
\label{desirables-section-proper-pushforward}

\noindent
代数スタックに対する Chow の補題が得られたので，現在はこれに取り組み始めることができる．
前節を参照せよ．


\section{Keel と Mori}
\label{desirables-section-keel-mori}

\noindent
\cite{K-M} を参照せよ．彼らの結果は，スタックの射の補足，節
\ref{stacks-more-morphisms-section-Keel-Mori} に追加されている．


\section{ここにさらに追加する}
\label{desirables-section-add-more}

\noindent
いや，そもそもこの一覧を Stacks プロジェクト自体の一部として始めるべきではなかった！
別の場所に，はるかに更新しやすい todo リストがある．
